\section{Reproduction, artifacts, and further work}
The release is self-contained for its executable checks. Python 3.10 or
later and the standard library suffice for the main replay. No network
access, solver, optimizer, Lean, or Rocq installation is needed.
From the unpacked top-level directory run
\begin{verbatim}
python3 -B reproduce.py
\end{verbatim}
The default command verifies the release hashes, binds the selected
parent and unchanged inference modules to the original Joint archive,
runs that parent's quick construction/source/startup replay in a temporary
directory, reconstructs both 208-state variants, checks the new local
theorem normally and under optimized Python, checks the alternative's
local histories, repeats the six-proposal comparison, verifies exact
startup, and rederives all 5,084 primary U208 contracts. Generated
temporary files do not alter the sealed release. To record a receipt,
pass \texttt{--report} with a path outside the release directory.

\begin{verbatim}
python3 -B reproduce.py --quick
\end{verbatim}
This command explicitly skips the fresh universal runtime derivation.
It remains useful for the new small proof and exact construction, but
must not be reported as a fresh full symbolic replay. The first complete
symbolic run in this preparation environment took about 139 seconds;
wall time is environment-dependent.

The optional independent finite executor requires a C++17 compiler:
\begin{verbatim}
g++ -std=c++17 -O3 code/check_finite_runtime.cpp -o /tmp/check_u208
/tmp/check_u208 machine/u208.native.tm machine/sites.tsv
\end{verbatim}
Rebuild the PDF in a working copy, so generated TeX auxiliary files do not
change the sealed inventory:
\begin{verbatim}
pdflatex -interaction=nonstopmode -halt-on-error report.tex
pdflatex -interaction=nonstopmode -halt-on-error report.tex
\end{verbatim}
PDF bytes may vary with the TeX version or metadata. Literal machine
and source bytes are deterministic.

\begin{center}
\small
\begin{tabular}{@{}>{\raggedright\arraybackslash}p{48mm}>{\raggedright\arraybackslash}p{100mm}@{}}
\toprule
Path & Contents\\
\midrule
\path{report.tex}, \path{report.pdf}, \path{paper/} & Editable report, proofs, comparison, full transition table\\
\path{machine/} & Primary native/dense U208 tables, state map, unchanged source\\
\path{alternative/} & Nonisomorphic Compatibility-derived U208 and its parent/source\\
\path{machine_parent/} & Exact selected U209 table and runtime label correspondence\\
\path{code/}, \path{reproduce.py} & Deterministic builder, independent invariant checker, runtime and startup replay\\
\path{comparison/} & All six table presentations, source snapshots, isomorphisms and fresh quick logs\\
\path{evidence/} & New local, startup, universal, finite, and release replay receipts\\
\path{reference/} & Full selected Joint archive with nested U212/U213/U214/U218/U222 lineage and unary theory; source provenance and reviews\\
\path{research/} & Discovery programs and precisely scoped search receipts\\
\path{licenses/}, \path{NOTICE.txt} & Upstream terms and attribution\\
\path{SHA256SUMS} & Hashes of all other sealed release files\\
\bottomrule
\end{tabular}
\end{center}

\subsection{What the new result suggests}
The effective improvement is a stronger invariant language for state
reuse. An unconditional incoming-bit test misses states entered from
different directions. Here a two-cell old-mode tag and a conditional
old-helper assertion suffice, even though the replacement path revisits
the same physical helper. Complete backward history gives discovery
candidates; small inductive assertions give a clearer proof.

A brief follow-up search on U208 found no further direct alias within
seven-step history contexts and 27 replacement steps, nor a further
shared-row fold across 348 ordered interfaces within the tested
seven-history/25-step limits. These are restricted negative searches,
not a lower bound at 208. Richer invariant combinations, several jointly
changed rows, or high-level arithmetic changes remain possible. The
report does not claim that any of those approaches already gives 207.

The immediate formalization target is compact: encode the five-state
invariant and four finite trace identities in the preferred arithmetic
or proof-assistant framework, then export the generic positive-block
simulation theorem for the exact table. Completing an end-to-end U208
certificate also requires importing or formalizing the retained
numerical-calculus and unary-theory proof transformations. That larger
task should be distinguished from the new finite reduction already
proved here.
